% CREATED BY DAVID FRISK, 2016
\chapter{Introduction}
\label{chap:introduction}

Modern vehicles generate large amounts of telemetry data through distributed embedded systems, sensors, actuators, and in-vehicle networks~\cite{Lee2017,Navet2009,MITRE2023}. These systems are designed to support real-time functionality, but they are also constrained by limited communication bandwidth, storage capacity, and on-board compute resources~\cite{Navet2009}. In practice, this has led to the widespread use of event-triggered and threshold-based logging strategies, which reduce data volume by transmitting only selected events or diagnostic messages~\cite{AUTOSAR2022b,Navet2009}. While these approaches are effective for limiting network load, they also risk omitting information that may later be relevant for offline analysis or machine learning-based applications~\cite{Navet2009,EventDrivenDiagPatent}.

This challenge has become more important as vehicles generate increasingly rich signal-based data and as downstream machine learning tasks have become a central use of automotive telemetry~\cite{Bertoncello2021,Samantaray2023,Theissler2021}. Collected vehicle data is no longer used only for diagnostics or monitoring, but also for tasks such as predictive maintenance, anomaly detection, fleet analytics, and other forms of offline model training~\cite{Theissler2021,Oezdemir2024,Brunheroto2022}. For these use cases, data utility matters as much as data volume: aggressive reduction of telemetry may improve efficiency, but it can also degrade the performance of models trained on the resulting data~\cite{Tirupathi2022,Eichinger2015,Moon2018}. This creates a fundamental tension between efficient storage and transmission on the one hand and preserving task-relevant information on the other.

The theoretical basis for this tension comes from information theory and rate-distortion theory~\cite{Shannon1948,Shannon1949}. Information entropy describes how much uncertainty is contained in a signal and therefore provides an indication of how compressible that signal is~\cite{Shannon1948}. Rate-distortion theory extends this idea to lossy compression by formalizing the trade-off between bitrate and reconstruction quality~\cite{Shannon1949}. For machine learning settings, this trade-off is not sufficient on its own, because good reconstruction does not necessarily imply that the compressed representation is useful for downstream tasks~\cite{Marchioni2024,Cuza2024}. This is why modern learned compression methods are often framed in terms of autoencoders and information bottlenecks: they aim to produce compact latent representations while retaining the information that is relevant for later prediction or classification tasks~\cite{Yang2022,Tishby2000,Chen2023,Alemi2016}.

Within this thesis, neural compression is the central approach for addressing this problem. Neural compression methods use learned encoder-decoder architectures to transform high-dimensional input signals into more compact representations. In contrast to traditional algorithmic compression methods, these models can adapt to the structure of the data and can be optimized for the kinds of signals that occur in automotive telemetry~\cite{Yang2022,Zheng2023,Liu2024}. In particular, vector quantization offers a practical way to introduce discrete bottlenecks that reduce entropy while still allowing the model to preserve useful structure in the compressed representation~\cite{Yang2022,Hodo2025}. This makes neural compression especially attractive for edge-oriented settings, where the compressed data must be stored, transmitted, or processed under strict efficiency constraints~\cite{Azar2022,Hodo2025}.

Recent research shows that neural compression has been successful in a range of domains, including time-series, audio, speech, and other resource-constrained applications~\cite{Löhdefink2019,Zheng2023,Liu2024,Zeghidour2021,Defossez2022,Ji2025,Azar2022,Hodo2025}. For time-series data, continuous-latent autoencoder approaches have demonstrated strong reconstruction performance, while vector-quantized methods such as EdgeCodec show that lightweight discrete compression can be effective in edge-deployment scenarios~\cite{Zheng2023,Hodo2025}. At the same time, the literature also shows that many methods primarily optimize reconstruction error, whereas downstream task utility is often treated as secondary or is not investigated in depth~\cite{Zheng2023,Hodo2025}. For vehicle telemetry, this is a crucial limitation, because a representation that reconstructs well is not necessarily the best representation for predictive maintenance or anomaly detection~\cite{Theissler2021,Oezdemir2024,Tirupathi2022}.

The background of this thesis therefore combines three strands of work: the properties and constraints of automotive telemetry, the information-theoretic foundations of compression, and the recent development of neural compression methods. Taken together, these areas motivate the central question of the thesis, namely how to compress multimodal vehicle telemetry in a way that remains useful for downstream machine learning tasks while still achieving meaningful compression and efficiency gains. The following sections build on this background by stating the concrete aim of the thesis, defining the research questions, and outlining the limitations of the study.

\section{Aim}
The aim of this thesis is threefold. First, two existing neural compression methods are applied to automotive and IoT time-series data in order to evaluate how they perform in terms of compression ratio, reconstruction quality, and computational efficiency. One of these methods is based on vector quantization, while the other serves as a continuous-latent baseline.

Second, the compressed representations produced by these methods are used to train downstream machine learning models and are compared against models trained on uncompressed data. This makes it possible to assess how much task-relevant information is retained after compression and reconstruction.

Third, the two methods are further explored through changes to the quantization strategy and loss function. The goal of this exploratory part is to investigate whether reconstruction loss and downstream utility retention can be improved further without sacrificing the general compression objective and computational efficiency.

\section{Research Questions}
This thesis is guided by the following research questions:

\begin{enumerate}
	\item RQ1: How do the selected neural compression methods perform on automotive and IoT time-series data in terms of compression ratio, reconstruction quality, and computational efficiency?
	\item RQ2: To what extent do downstream machine learning models trained on reconstructed compressed data retain task performance compared to models trained on uncompressed data?
	\item RQ3: Where do the existing neural compression methods fall short in terms of reconstruction quality and downstream utility retention, and which quantization strategies and loss functions can mitigate these limitations?
\end{enumerate}

\section{Limitations}
Due to time constraints and the scope of this thesis, there are several limitations to the work presented. 

First, the work forgoes testing on real hardware as that is very time consuming and difficult to set up while contributing relatively little to the main research question. Instead, canonical methods to measure computational efficiency are employed. 

The work is further limited to testing only regression tasks and anomaly detection as downstream tasks due to time constraints and the structure of the available datasets. These two tasks should, however, be a good representation of common downstream tasks, and adding more tasks would likely not drastically change the conclusion of the thesis. 

To keep the compressed data representation as versatile as possible, no task-aware compression techniques are implemented. While task-aware compression could potentially yield better performance for specific tasks, it would limit the generalizability of the compressed representation and add significant complexity to the work. It also results in a more relevant comparison to other compression techniques which are not task-aware. 


% This chapter presents the section levels that can be used in the template. 

% \section{Section levels}
% \autoref{tab:sections} presents an overview of the section levels that are used in this document. The number of levels that are numbered and included in the table of contents is set in the settings file \texttt{Settings.tex}. The levels are shown in Section \ref{Section_ref}.

% This is a new paragraph and should have proper parskip or indentation. Don't forget to cite your sources~\cite{Brajnik2008}. % '~' becomes space which cannot line break.

% \begin{table}[h]
% \centering
% \caption{Section levels} % Table text above table.
% \begin{tabular}{ll} \hline
% Name & Command\\ \hline
% Chapter & \textbackslash\texttt{chapter\{\emph{Chapter name}\}}\\
% Section & \textbackslash\texttt{section\{\emph{Section name}\}}\\
% Subsection & \textbackslash\texttt{subsection\{\emph{Subsection name}\}}\\
% Subsubsection & \textbackslash\texttt{subsubsection\{\emph{Subsubsection name}\}}\\
% %Paragraph & \textbackslash\texttt{paragraph\{\emph{Paragraph name}\}}\\
% %Subparagraph & \textbackslash\texttt{paragraph\{\emph{Subparagraph name}\}}\\ \hline\hline
% \end{tabular}
% \label{tab:sections}
% \end{table}

% \section{Section} \label{Section_ref}
% \subsection{Subsection}
% \subsubsection{Subsubsection}
% %\paragraph{Paragraph}
% %\subparagraph{Subparagraph}

